\section{Discussion}\label{sec:discussion}

\subsection{When does distribution pay off?}
Every cluster run pays a fixed cost that does not depend on the input: about 13--18\,s
for Spark and 18--23\,s for MapReduce on the local cluster (the tier~S time minus the
serial time), and about twice as much on the slower cloud VMs. On one month of data
this cost is three to five times the whole serial computation, so a cluster is
pointless at that size. Spark DataFrame on one worker breaks even with the serial
baseline at about six months of data (tier M, 1\,GB) and is 1.2--1.4$\times$ faster at
twelve months. MapReduce needs three to four workers to match the serial program on 2\,GB and wins clearly only
on 3.7\,GB (tier XL). Note that one worker has two CPUs while the serial program uses
one, so this comparison even favours the cluster.

The weak-scaling efficiency of 0.67--0.79 at four workers shows the same effect from
the other side: when the data grow with the cluster, the run time grows only slowly,
mainly because the fixed cost is paid once per job rather than once per worker.

\subsection{Scalability and its limits}
\textbf{Amdahl's law.} The Karp--Flatt estimate on tier L at $N=4$ is
$e=0.15$--$0.22$ for MapReduce, 0.23--0.33 for Spark RDD and 0.37--0.44 for Spark
DataFrame (\cref{tab:summary}). These values roughly agree with a simple model in
which the fixed cost is the serial part: for Spark DataFrame, about 15\,s of fixed cost
out of $T_1\approx50$\,s is a serial fraction of about 0.3, which by Amdahl's law
bounds the speedup to about~3.3 however many workers are added; for MapReduce, about
20\,s out of 150\,s gives 0.13. The fastest engine therefore scales
\emph{worst}: its parallel part is so short that the constant start-up dominates. On
tier XL the parallel part is larger and $e$ drops to 0.10--0.26 for Spark DataFrame and
below 0.07 for MapReduce. Relative speedup alone would make MapReduce look like the
better system; the absolute times show the opposite.

\textbf{Task granularity.} Tier L has 24 HDFS blocks and hence 24 map tasks. YARN
schedules by memory, so each worker runs three 1\,GB containers, one of which hosts
the MapReduce application master. With three workers 8 maps run at once (exactly three
waves); with four workers 11 run at once, and the third wave contains only two tasks.
This explains why MapReduce gains only 4--15\% from the fourth worker on tier L but
3.35--3.97$\times$ from one to four workers on tier XL (48 tasks).

\begin{table}[t]
\centering
\caption{Cache experiment, \texttt{routes} on tier L: mean run time (s) after dropping
the page cache (cold) and directly afterwards (warm), largest disk read in a warm run,
and CPU time withheld by the 2-CPU cap per worker and second (3 pairs per cell)}\label{tab:cache}
\setlength{\tabcolsep}{2.5pt}
\begin{tabular*}{\columnwidth}{@{\extracolsep{\fill}}rrrrrrrrr@{}}
\toprule
 & Data/wkr & \multicolumn{2}{c}{MapReduce} & \multicolumn{2}{c}{Spark DF} & Warm read & \multicolumn{2}{c}{Throttled} \\
\cmidrule(lr){3-4}\cmidrule(lr){5-6}\cmidrule(l){8-9}
$N$ & (GiB) & cold & warm & cold & warm & (MiB) & MR & DF \\
\midrule
\inputrows{generated/cache_table.tex}
\bottomrule
\end{tabular*}
\end{table}

\textbf{The two-worker anomaly.} The poor result at $N=2$ on tier L
(\cref{sec:results}) is specific to one tier and one environment. Our first
hypothesis was the page cache: the 5\,GB limit of a container also counts the cached
HDFS blocks, so with one or two workers holding 1--2\,GB each, blocks might be re-read
from the shared virtual disk. We tested it with a further 48 runs
(\cref{tab:cache}): \texttt{routes} on tier L with $N=1$--4, MapReduce and Spark
DataFrame, three pairs each, where the page cache of the Docker VM is dropped before
the \emph{cold} run and the \emph{warm} run follows directly. For each worker we read
the cgroup counters for bytes read from disk, refaulted cache pages and throttled CPU
time. The hypothesis is \emph{refuted}: even with all 2.3\,GiB on one worker, warm
runs read at most 3\,MiB from disk, no page was refaulted and no container reached its
memory limit, and cold runs, which read the whole input from disk, differ from warm
runs by less than 5\,s. Moreover, the anomaly did not reproduce: at $N=2$ the rerun gives
$S_\text{rel}(2)\approx1.5$ for both engines, on a smooth curve between $N=1$ and
$N=3$. In the main benchmark all $N=2$ runs on tier L fell into one hour of the night
(01:30--02:30) and the configurations ran in a fixed order, so we attribute the anomaly
to a transient disturbance on the host during that hour, which we cannot identify
after the fact. The rerun also shows the size of session-to-session variation on a
shared host: at $N=1$ it matches the main benchmark within 3\%, but at $N=3$--4 its
warm runs are 20\% slower, so MapReduce reaches only $S_\text{rel}(4)\approx2.0$ instead of 2.4.
Speedups on tier L therefore carry an uncertainty of this order.

The throttling counters support the task-granularity argument. A MapReduce worker runs
three 1\,GB containers, each a JVM plus a Python process, on a 2-CPU quota: the cap
withholds 0.7 CPU-seconds per second at $N=1$, where one container is the application
master, and 1.4--1.7 at $N\ge2$, where a worker runs three map tasks. Spark runs one
2-core executor per worker and is throttled by only 0.2. MapReduce is thus CPU-bound
and oversubscribed. Finally, the \texttt{cpus} limit caps CPU time but does not pin
containers to cores, so on the hybrid i5-13600KF a worker may run on a performance core,
an efficiency core or a hyper-thread of a busy core.

\subsection{Effect of the combiner}
The combiner reduces the shuffle by two to five orders of magnitude
(\cref{fig:combcloud}a), but saves time only when the shuffle is large compared with
the rest of the job. On \texttt{fare} the uncombined shuffle (1.1\,GiB) is 2.5--2.8
times that of \texttt{demand} or \texttt{routes}, and the combiner saves up to 24\%
locally and 29\% on Azure. For the other two jobs, 400--450\,MiB are shuffled without measurable
cost on one host, where the ``network'' is a virtual bridge in memory, but the combiner
saves 7--10\% on Azure, where the data cross a real network. On a larger cluster with
slower links we expect the benefit to grow.

\subsection{Hadoop MapReduce vs.\ Spark}
On one worker and tier L, Spark RDD is 2.0--2.1$\times$ faster than MapReduce and Spark
DataFrame 3.0--3.1$\times$ faster, although all three run the same algorithm. MapReduce
starts a new JVM for every task, writes and sorts all map output on local disk before
the reduce phase, and, with Streaming, passes every record as text through a pipe to a
Python process and back. Spark reuses one executor per worker for all tasks and
aggregates in an in-memory hash table~\cite{zaharia2012rdd}. Spark RDD still calls the
Python \texttt{parse} and \texttt{map\_record} functions once per record and moves data
between the JVM and Python workers; the DataFrame version runs entirely in the JVM.
The 1.5$\times$ difference between Spark RDD and DataFrame is therefore the cost of
executing user code in Python, and the remaining 2$\times$ between MapReduce and Spark
RDD is the cost of the disk-based MapReduce execution model.

\subsection{Single host vs.\ public cloud}
The local machine is faster per worker: the serial program runs 2.0$\times$ faster on
one i5-13600KF core than on one Azure vCPU, and the engines 2.5--3.3$\times$ faster.
Two Azure vCPUs are the two hyper-threads of one older server core, whereas a local
worker usually gets two separate modern cores. The cloud, however, scales more
regularly, because each worker there owns its disk, memory and network interface. The
single-host cluster is cheap and convenient for development, but adding containers on
one host adds CPU quota, not independent hardware, and its speedup measurements can be
distorted by shared resources. The cloud cluster cost about US\$4 for all experiments;
for the same money the local machine was faster, but it cannot grow beyond one host.

\subsection{Lessons learned}
(1)~\emph{Measure the fixed cost first.} One run on one month told us more about when
a cluster pays off than any estimate: below about 1\,GB the start-up cost dominates.
(2)~\emph{Containers are not machines.} Without caps one container used the whole host;
even with caps they share a disk, and our measurements varied by 20\% between
sessions, so multi-machine claims need real machines. (3)~\emph{Write the logic once,
check every output.} One shared parser and an automatic output comparison made every
speed comparison fair. (4)~\emph{Automate and make runs resumable.} The 10-hour sweep
was interrupted several times; recording each run let us resume instead of restarting.
(5)~\emph{Test an explanation before writing it.} Our first explanation of the
two-worker anomaly, the page cache, was wrong.

\subsection{Threats to validity}
(i)~All local workers share one host, disk and memory bus (see above). (ii)~Only two
cloud workers were possible within the student subscription's limit of six vCPUs per
region, and the cloud master uses a different CPU model than the workers.
(iii)~Repetitions were consecutive and the configurations ran in a fixed order, so a
disturbance of the host affects all repetitions of one configuration, as at $N=2$ on
tier L; a rerun in another session differed by up to 20\%. A warm page cache, in
contrast, changes run times by at most 5\,s (\cref{tab:cache}). (iv)~Apart from the
cache experiment we measured wall-clock time only, and even there not network
traffic or per-core CPU utilisation. (v)~The input is CSV converted from Parquet,
which favours line-based Hadoop Streaming; Spark could read Parquet directly.
